% Options for packages loaded elsewhere
\PassOptionsToPackage{unicode}{hyperref}
\PassOptionsToPackage{hyphens}{url}
%
\documentclass[
]{article}
\usepackage{amsmath,amssymb}
\usepackage{iftex}
\ifPDFTeX
  \usepackage[T1]{fontenc}
  \usepackage[utf8]{inputenc}
  \usepackage{textcomp} % provide euro and other symbols
\else % if luatex or xetex
  \usepackage{unicode-math} % this also loads fontspec
  \defaultfontfeatures{Scale=MatchLowercase}
  \defaultfontfeatures[\rmfamily]{Ligatures=TeX,Scale=1}
\fi
\usepackage{lmodern}
\ifPDFTeX\else
  % xetex/luatex font selection
\fi
% Use upquote if available, for straight quotes in verbatim environments
\IfFileExists{upquote.sty}{\usepackage{upquote}}{}
\IfFileExists{microtype.sty}{% use microtype if available
  \usepackage[]{microtype}
  \UseMicrotypeSet[protrusion]{basicmath} % disable protrusion for tt fonts
}{}
\makeatletter
\@ifundefined{KOMAClassName}{% if non-KOMA class
  \IfFileExists{parskip.sty}{%
    \usepackage{parskip}
  }{% else
    \setlength{\parindent}{0pt}
    \setlength{\parskip}{6pt plus 2pt minus 1pt}}
}{% if KOMA class
  \KOMAoptions{parskip=half}}
\makeatother
\usepackage{xcolor}
\usepackage[margin=1in]{geometry}
\usepackage{color}
\usepackage{fancyvrb}
\newcommand{\VerbBar}{|}
\newcommand{\VERB}{\Verb[commandchars=\\\{\}]}
\DefineVerbatimEnvironment{Highlighting}{Verbatim}{commandchars=\\\{\}}
% Add ',fontsize=\small' for more characters per line
\usepackage{framed}
\definecolor{shadecolor}{RGB}{248,248,248}
\newenvironment{Shaded}{\begin{snugshade}}{\end{snugshade}}
\newcommand{\AlertTok}[1]{\textcolor[rgb]{0.94,0.16,0.16}{#1}}
\newcommand{\AnnotationTok}[1]{\textcolor[rgb]{0.56,0.35,0.01}{\textbf{\textit{#1}}}}
\newcommand{\AttributeTok}[1]{\textcolor[rgb]{0.13,0.29,0.53}{#1}}
\newcommand{\BaseNTok}[1]{\textcolor[rgb]{0.00,0.00,0.81}{#1}}
\newcommand{\BuiltInTok}[1]{#1}
\newcommand{\CharTok}[1]{\textcolor[rgb]{0.31,0.60,0.02}{#1}}
\newcommand{\CommentTok}[1]{\textcolor[rgb]{0.56,0.35,0.01}{\textit{#1}}}
\newcommand{\CommentVarTok}[1]{\textcolor[rgb]{0.56,0.35,0.01}{\textbf{\textit{#1}}}}
\newcommand{\ConstantTok}[1]{\textcolor[rgb]{0.56,0.35,0.01}{#1}}
\newcommand{\ControlFlowTok}[1]{\textcolor[rgb]{0.13,0.29,0.53}{\textbf{#1}}}
\newcommand{\DataTypeTok}[1]{\textcolor[rgb]{0.13,0.29,0.53}{#1}}
\newcommand{\DecValTok}[1]{\textcolor[rgb]{0.00,0.00,0.81}{#1}}
\newcommand{\DocumentationTok}[1]{\textcolor[rgb]{0.56,0.35,0.01}{\textbf{\textit{#1}}}}
\newcommand{\ErrorTok}[1]{\textcolor[rgb]{0.64,0.00,0.00}{\textbf{#1}}}
\newcommand{\ExtensionTok}[1]{#1}
\newcommand{\FloatTok}[1]{\textcolor[rgb]{0.00,0.00,0.81}{#1}}
\newcommand{\FunctionTok}[1]{\textcolor[rgb]{0.13,0.29,0.53}{\textbf{#1}}}
\newcommand{\ImportTok}[1]{#1}
\newcommand{\InformationTok}[1]{\textcolor[rgb]{0.56,0.35,0.01}{\textbf{\textit{#1}}}}
\newcommand{\KeywordTok}[1]{\textcolor[rgb]{0.13,0.29,0.53}{\textbf{#1}}}
\newcommand{\NormalTok}[1]{#1}
\newcommand{\OperatorTok}[1]{\textcolor[rgb]{0.81,0.36,0.00}{\textbf{#1}}}
\newcommand{\OtherTok}[1]{\textcolor[rgb]{0.56,0.35,0.01}{#1}}
\newcommand{\PreprocessorTok}[1]{\textcolor[rgb]{0.56,0.35,0.01}{\textit{#1}}}
\newcommand{\RegionMarkerTok}[1]{#1}
\newcommand{\SpecialCharTok}[1]{\textcolor[rgb]{0.81,0.36,0.00}{\textbf{#1}}}
\newcommand{\SpecialStringTok}[1]{\textcolor[rgb]{0.31,0.60,0.02}{#1}}
\newcommand{\StringTok}[1]{\textcolor[rgb]{0.31,0.60,0.02}{#1}}
\newcommand{\VariableTok}[1]{\textcolor[rgb]{0.00,0.00,0.00}{#1}}
\newcommand{\VerbatimStringTok}[1]{\textcolor[rgb]{0.31,0.60,0.02}{#1}}
\newcommand{\WarningTok}[1]{\textcolor[rgb]{0.56,0.35,0.01}{\textbf{\textit{#1}}}}
\usepackage{longtable,booktabs,array}
\usepackage{calc} % for calculating minipage widths
% Correct order of tables after \paragraph or \subparagraph
\usepackage{etoolbox}
\makeatletter
\patchcmd\longtable{\par}{\if@noskipsec\mbox{}\fi\par}{}{}
\makeatother
% Allow footnotes in longtable head/foot
\IfFileExists{footnotehyper.sty}{\usepackage{footnotehyper}}{\usepackage{footnote}}
\makesavenoteenv{longtable}
\usepackage{graphicx}
\makeatletter
\newsavebox\pandoc@box
\newcommand*\pandocbounded[1]{% scales image to fit in text height/width
  \sbox\pandoc@box{#1}%
  \Gscale@div\@tempa{\textheight}{\dimexpr\ht\pandoc@box+\dp\pandoc@box\relax}%
  \Gscale@div\@tempb{\linewidth}{\wd\pandoc@box}%
  \ifdim\@tempb\p@<\@tempa\p@\let\@tempa\@tempb\fi% select the smaller of both
  \ifdim\@tempa\p@<\p@\scalebox{\@tempa}{\usebox\pandoc@box}%
  \else\usebox{\pandoc@box}%
  \fi%
}
% Set default figure placement to htbp
\def\fps@figure{htbp}
\makeatother
\setlength{\emergencystretch}{3em} % prevent overfull lines
\providecommand{\tightlist}{%
  \setlength{\itemsep}{0pt}\setlength{\parskip}{0pt}}
\setcounter{secnumdepth}{-\maxdimen} % remove section numbering
\usepackage{bookmark}
\IfFileExists{xurl.sty}{\usepackage{xurl}}{} % add URL line breaks if available
\urlstyle{same}
\hypersetup{
  pdftitle={Examen visualisation: International Trade Network},
  hidelinks,
  pdfcreator={LaTeX via pandoc}}

\title{Examen visualisation: International Trade Network}
\author{}
\date{\vspace{-2.5em}19 juillet 2025}

\begin{document}
\maketitle

La taille et la structure des flux commerciaux internationaux varient
considérablement au fil du temps. Cet exercice s'appuie en partie sur
Luca De Benedictis et Lucia Tajoli. (2011). `The World Trade Network.'
\emph{The World Economy}, 34:8, pp.1417-1454. Les données commerciales
proviennent de Katherine Barbieri et Omar Keshk. (2012).
\emph{Correlates of War Project Trade Data Set}, Version 3.0. et sont
disponibles à : \url{http://correlatesofwar.org}.

Le volume des biens échangés entre les pays a augmenté rapidement au
cours du siècle dernier, à mesure que les progrès technologiques ont
réduit le coût du transport et que les pays ont adopté des politiques
commerciales plus libérales. Parfois, cependant, les flux commerciaux
ont diminué en raison d'événements perturbateurs tels que des guerres
majeures et l'adoption de politiques commerciales protectionnistes. Dans
cet exercice, nous allons explorer certains de ces changements en
examinant le réseau de commerce international sur plusieurs périodes. Le
fichier de données \texttt{trade.csv} contient la valeur des
exportations d'un pays vers un autre au cours d'une année donnée. Les
noms et descriptions des variables de cet ensemble de données sont :

\begin{longtable}[]{@{}ll@{}}
\toprule\noalign{}
Name & Description \\
\midrule\noalign{}
\endhead
\bottomrule\noalign{}
\endlastfoot
\texttt{country1} & Country name of exporter \\
\texttt{country2} & Country name of importer \\
\texttt{year} & Year \\
\texttt{exports} & Total value of exports (in tens of millions of
dollars) \\
\end{longtable}

Les données sont données pour les années 1900, 1920, 1940, 1955, 1980,
2000 et 2009.

\subsection{Question 1}\label{question-1}

Après avoir lu l'article, faites en un résumé succinct des objectifs et
de la méthodologie. décrivez les principaux indicateurs de mesure des
données de réseaux

REPONSE: L'objectif principal de l'article est d'analyser le commerce
international à travers une approche de réseau. Au lieu de voir les
échanges entre pays comme des données tabulaires classiques, les auteurs
les modélisent comme un réseau dirigé où :

Chaque pays est un nœud,

Chaque flux commercial est un arc allant d'un pays exportateur vers un
pays importateur.

Cette approche permet de mieux comprendre la structure globale du
système commercial mondial, ses dynamiques historiques (de 1950 à 2000)
et les positions stratégiques des pays dans le réseau international.

Méthodologie Modélisation du réseau : Les auteurs construisent pour
chaque année un graphe dirigé, pondéré ou non, selon que l'on tient
compte de la valeur des exportations ou uniquement de leur existence (0
ou 1).

Construction des matrices d'adjacence : Ils créent des matrices où
chaque cellule (i,j) représente un lien d'exportation du pays i vers le
pays j.

Analyse temporelle : L'étude est menée pour plusieurs années clefs
(notamment 1950, 1980 et 2000), pour observer les changements
structurels du commerce mondial.

Mesures de réseau : Divers indicateurs issus de la théorie des graphes
sont utilisés pour quantifier la structure du réseau et le rôle des
pays.

Principaux indicateurs de réseau utilisés Densité du réseau (network
density) :

Proportion de liens effectivement existants par rapport au nombre
maximal possible.

Donne une idée du niveau global de connectivité du réseau.

Degré (degree centrality) :

Nombre de partenaires commerciaux d'un pays.

On distingue le degré entrant (nombre d'importateurs) et le degré
sortant (nombre d'exportateurs).

Centralité d'intermédiarité (betweenness centrality) :

Mesure combien de fois un pays agit comme pont entre d'autres pays dans
les chemins commerciaux les plus courts.

Indique les pays-clés dans la circulation mondiale des biens.

Proximité (closeness centrality) :

Mesure l'accessibilité d'un pays aux autres dans le réseau.

Plus un pays est « proche » de tous les autres, plus il peut rapidement
échanger.

Composantes fortement connexes :

Sous-groupes de pays qui échangent entre eux de manière réciproque.

Permet d'identifier des blocs régionaux ou géopolitiques (ex : Europe de
l'Ouest, ex-bloc soviétique).

Conclusion

L'article applique une approche de réseau pour analyser la structure du
commerce international entre pays, en mesurant la densité, la centralité
(degré, intermédiarité, proximité) et les composantes du réseau
commercial mondial sur plus d'un demi-siècle.

\subsection{Question 2}\label{question-2}

On n'a pas les mêmes pays chaque année. Certains pays disparaissent
alors que d'autres naissent, surtout lors de la dislocation du bloc
soviétique. Créer une list qui fait la liste des pays intervenus dans le
commerce chaque année

\begin{Shaded}
\begin{Highlighting}[]
\CommentTok{\# Charger les données}

\FunctionTok{library}\NormalTok{(readr)}
\NormalTok{trade }\OtherTok{\textless{}{-}} \FunctionTok{read\_csv}\NormalTok{(}\StringTok{"data/trade.csv"}\NormalTok{)}
\end{Highlighting}
\end{Shaded}

\begin{verbatim}
## Rows: 114588 Columns: 4
## -- Column specification --------------------------------------------------------
## Delimiter: ","
## chr (2): country1, country2
## dbl (2): year, exports
## 
## i Use `spec()` to retrieve the full column specification for this data.
## i Specify the column types or set `show_col_types = FALSE` to quiet this message.
\end{verbatim}

\begin{Shaded}
\begin{Highlighting}[]
\CommentTok{\# Remplacer les NA dans la variable exports par 0}
\NormalTok{trade}\SpecialCharTok{$}\NormalTok{exports[}\FunctionTok{is.na}\NormalTok{(trade}\SpecialCharTok{$}\NormalTok{exports)] }\OtherTok{\textless{}{-}} \DecValTok{0}


\CommentTok{\# Créer une liste des pays impliqués dans le commerce pour chaque année}
\FunctionTok{library}\NormalTok{(dplyr)}
\end{Highlighting}
\end{Shaded}

\begin{verbatim}
## 
## Attachement du package : 'dplyr'
## 
## Les objets suivants sont masqués depuis 'package:stats':
## 
##     filter, lag
## 
## Les objets suivants sont masqués depuis 'package:base':
## 
##     intersect, setdiff, setequal, union
\end{verbatim}

\begin{Shaded}
\begin{Highlighting}[]
\NormalTok{pays\_par\_annee }\OtherTok{\textless{}{-}}\NormalTok{ trade }\SpecialCharTok{\%\textgreater{}\%}
  \FunctionTok{group\_by}\NormalTok{(year) }\SpecialCharTok{\%\textgreater{}\%}
  \FunctionTok{summarise}\NormalTok{(}
    \AttributeTok{pays =} \FunctionTok{list}\NormalTok{(}\FunctionTok{sort}\NormalTok{(}\FunctionTok{unique}\NormalTok{(}\FunctionTok{c}\NormalTok{(country1, country2)))),}
    \AttributeTok{nb\_pays =} \FunctionTok{length}\NormalTok{(}\FunctionTok{unique}\NormalTok{(}\FunctionTok{c}\NormalTok{(country1, country2)))}
\NormalTok{  )}



\FunctionTok{library}\NormalTok{(openxlsx)}

\CommentTok{\# Exporter le tableau en Excel}
\FunctionTok{write.xlsx}\NormalTok{(pays\_par\_annee, }\StringTok{"data/pays\_par\_annee.xlsx"}\NormalTok{)}

\CommentTok{\# Afficher le résultat}
\FunctionTok{print}\NormalTok{(pays\_par\_annee)}
\end{Highlighting}
\end{Shaded}

\begin{verbatim}
## # A tibble: 7 x 3
##    year pays        nb_pays
##   <dbl> <list>        <int>
## 1  1900 <chr [42]>       42
## 2  1920 <chr [59]>       59
## 3  1940 <chr [62]>       62
## 4  1955 <chr [84]>       84
## 5  1980 <chr [156]>     156
## 6  2000 <chr [191]>     191
## 7  2009 <chr [196]>     196
\end{verbatim}

\subsection{Question 3}\label{question-3}

Maintenant, on va regarder uniquement les données de 1900. Créer une
nouvelle base de données uniquement avec cette année. Quelle est la
somme totale des exportations des états unis cette année. Quelle est la
part de cette exportation vers tous les pays avec lesquels les États
unis ont commercé cette année? Présenter un graphique qui résume votre
réponse à la dernière question.

\begin{Shaded}
\begin{Highlighting}[]
\CommentTok{\# 1. Filtrer les données pour l\textquotesingle{}année 1900}
\NormalTok{data\_1900 }\OtherTok{\textless{}{-}}\NormalTok{ trade }\SpecialCharTok{\%\textgreater{}\%} \FunctionTok{filter}\NormalTok{(year }\SpecialCharTok{==} \DecValTok{1900}\NormalTok{)}

\CommentTok{\# 2. Extraire les exportations depuis les États{-}Unis}
\NormalTok{usa\_exports\_1900 }\OtherTok{\textless{}{-}}\NormalTok{ data\_1900 }\SpecialCharTok{\%\textgreater{}\%} 
  \FunctionTok{filter}\NormalTok{(country1 }\SpecialCharTok{==} \StringTok{"United States of America"}\NormalTok{)}

\CommentTok{\# 3. Calculer la somme totale des exportations américaines en 1900}
\NormalTok{total\_exports\_usa }\OtherTok{\textless{}{-}} \FunctionTok{sum}\NormalTok{(usa\_exports\_1900}\SpecialCharTok{$}\NormalTok{exports, }\AttributeTok{na.rm =} \ConstantTok{TRUE}\NormalTok{)}

\CommentTok{\# 4. Calculer la part de chaque pays importateur dans les exportations US}
\NormalTok{usa\_exports\_1900 }\OtherTok{\textless{}{-}}\NormalTok{ usa\_exports\_1900 }\SpecialCharTok{\%\textgreater{}\%}
  \FunctionTok{mutate}\NormalTok{(}
    \AttributeTok{part =}\NormalTok{ exports }\SpecialCharTok{/}\NormalTok{ total\_exports\_usa  }\CommentTok{\# proportion des exportations par pays}
\NormalTok{  )}

\CommentTok{\# 5. Graphique : part des exportations par pays importateur (barplot)}
\FunctionTok{library}\NormalTok{(ggplot2)}

\FunctionTok{ggplot}\NormalTok{(usa\_exports\_1900, }\FunctionTok{aes}\NormalTok{(}\AttributeTok{x =} \FunctionTok{reorder}\NormalTok{(country2, }\SpecialCharTok{{-}}\NormalTok{part), }\AttributeTok{y =}\NormalTok{ part)) }\SpecialCharTok{+}
  \FunctionTok{geom\_bar}\NormalTok{(}\AttributeTok{stat =} \StringTok{"identity"}\NormalTok{, }\AttributeTok{fill =} \StringTok{"steelblue"}\NormalTok{) }\SpecialCharTok{+}
  \FunctionTok{labs}\NormalTok{(}
    \AttributeTok{title =} \StringTok{"Part des exportations américaines par pays en 1900"}\NormalTok{,}
    \AttributeTok{x =} \StringTok{"Pays importateur"}\NormalTok{,}
    \AttributeTok{y =} \StringTok{"Part (\%)"}
\NormalTok{  ) }\SpecialCharTok{+}
  \FunctionTok{scale\_y\_continuous}\NormalTok{(}\AttributeTok{labels =}\NormalTok{ scales}\SpecialCharTok{::}\NormalTok{percent) }\SpecialCharTok{+}
  \FunctionTok{theme\_minimal}\NormalTok{() }\SpecialCharTok{+}
  \FunctionTok{theme}\NormalTok{(}\AttributeTok{axis.text.x =} \FunctionTok{element\_text}\NormalTok{(}\AttributeTok{angle =} \DecValTok{60}\NormalTok{, }\AttributeTok{hjust =} \DecValTok{1}\NormalTok{))}
\end{Highlighting}
\end{Shaded}

\pandocbounded{\includegraphics[keepaspectratio]{Examen_visualisation_trade-networks_files/figure-latex/unnamed-chunk-2-1.pdf}}

\begin{Shaded}
\begin{Highlighting}[]
\CommentTok{\# Garder uniquement les 10 premiers pays par valeur d’exportation}
\NormalTok{top10\_usa\_exports }\OtherTok{\textless{}{-}}\NormalTok{ usa\_exports\_1900 }\SpecialCharTok{\%\textgreater{}\%}
  \FunctionTok{arrange}\NormalTok{(}\FunctionTok{desc}\NormalTok{(exports)) }\SpecialCharTok{\%\textgreater{}\%}
  \FunctionTok{slice}\NormalTok{(}\DecValTok{1}\SpecialCharTok{:}\DecValTok{10}\NormalTok{)}

\CommentTok{\# Graphique plus lisible (Top 10 uniquement)}
\FunctionTok{ggplot}\NormalTok{(top10\_usa\_exports, }\FunctionTok{aes}\NormalTok{(}\AttributeTok{x =} \FunctionTok{reorder}\NormalTok{(country2, }\SpecialCharTok{{-}}\NormalTok{part), }\AttributeTok{y =}\NormalTok{ part)) }\SpecialCharTok{+}
  \FunctionTok{geom\_bar}\NormalTok{(}\AttributeTok{stat =} \StringTok{"identity"}\NormalTok{, }\AttributeTok{fill =} \StringTok{"darkgreen"}\NormalTok{) }\SpecialCharTok{+}
  \FunctionTok{labs}\NormalTok{(}
    \AttributeTok{title =} \StringTok{"Top 10 des partenaires commerciaux des USA en 1900"}\NormalTok{,}
    \AttributeTok{x =} \StringTok{"Pays importateur"}\NormalTok{,}
    \AttributeTok{y =} \StringTok{"Part (\%)"}
\NormalTok{  ) }\SpecialCharTok{+}
  \FunctionTok{scale\_y\_continuous}\NormalTok{(}\AttributeTok{labels =}\NormalTok{ scales}\SpecialCharTok{::}\FunctionTok{percent\_format}\NormalTok{(}\AttributeTok{accuracy =} \DecValTok{1}\NormalTok{)) }\SpecialCharTok{+}
  \FunctionTok{theme\_minimal}\NormalTok{() }\SpecialCharTok{+}
  \FunctionTok{theme}\NormalTok{(}\AttributeTok{axis.text.x =} \FunctionTok{element\_text}\NormalTok{(}\AttributeTok{angle =} \DecValTok{60}\NormalTok{, }\AttributeTok{hjust =} \DecValTok{1}\NormalTok{))}
\end{Highlighting}
\end{Shaded}

\pandocbounded{\includegraphics[keepaspectratio]{Examen_visualisation_trade-networks_files/figure-latex/unnamed-chunk-3-1.pdf}}

\subsection{Question 4}\label{question-4}

Une meilleure représentation serait de faire le treemap. Présenter
l'information avec ce graphique en indiquant clairement les pays et les
continents en jeu.

\begin{Shaded}
\begin{Highlighting}[]
\CommentTok{\# Charger les bibliothèques nécessaires}
\FunctionTok{library}\NormalTok{(dplyr)}
\FunctionTok{library}\NormalTok{(ggplot2)}
\FunctionTok{library}\NormalTok{(treemapify)}
\FunctionTok{library}\NormalTok{(countrycode)}

\CommentTok{\# 1. Reprendre les exportations des USA en 1900 (déjà filtrées précédemment)}
\NormalTok{usa\_exports\_1900 }\OtherTok{\textless{}{-}}\NormalTok{ trade }\SpecialCharTok{\%\textgreater{}\%} \FunctionTok{filter}\NormalTok{(year }\SpecialCharTok{==} \DecValTok{1900}\NormalTok{, country1 }\SpecialCharTok{==} \StringTok{"United States of America"}\NormalTok{)}

\CommentTok{\# Correction manuelle des noms historiques}
\NormalTok{usa\_exports\_1900 }\OtherTok{\textless{}{-}}\NormalTok{ usa\_exports\_1900 }\SpecialCharTok{\%\textgreater{}\%}
  \FunctionTok{mutate}\NormalTok{(}
    \AttributeTok{country2 =} \FunctionTok{case\_when}\NormalTok{(}
\NormalTok{      country2 }\SpecialCharTok{==} \StringTok{"Austria{-}Hungary"} \SpecialCharTok{\textasciitilde{}} \StringTok{"Austria"}\NormalTok{,       }\CommentTok{\# ou "Hungary" selon ton choix}
\NormalTok{      country2 }\SpecialCharTok{==} \StringTok{"Yugoslavia"} \SpecialCharTok{\textasciitilde{}} \StringTok{"Serbia"}\NormalTok{,             }\CommentTok{\# ou "Bosnia", "Croatia", etc.}
      \ConstantTok{TRUE} \SpecialCharTok{\textasciitilde{}}\NormalTok{ country2}
\NormalTok{    ),}
    \AttributeTok{continent =} \FunctionTok{countrycode}\NormalTok{(country2, }\AttributeTok{origin =} \StringTok{"country.name"}\NormalTok{, }\AttributeTok{destination =} \StringTok{"continent"}\NormalTok{)}
\NormalTok{  )}


\CommentTok{\# 2. Ajouter une colonne "continent" basée sur les noms de pays importateurs}
\NormalTok{usa\_exports\_1900 }\OtherTok{\textless{}{-}}\NormalTok{ usa\_exports\_1900 }\SpecialCharTok{\%\textgreater{}\%}
  \FunctionTok{mutate}\NormalTok{(}
    \AttributeTok{continent =} \FunctionTok{countrycode}\NormalTok{(country2, }\AttributeTok{origin =} \StringTok{"country.name"}\NormalTok{, }\AttributeTok{destination =} \StringTok{"continent"}\NormalTok{)}
\NormalTok{  )}

\CommentTok{\# 3. Supprimer les lignes sans continent identifié}
\NormalTok{usa\_exports\_1900 }\OtherTok{\textless{}{-}}\NormalTok{ usa\_exports\_1900 }\SpecialCharTok{\%\textgreater{}\%} \FunctionTok{filter}\NormalTok{(}\SpecialCharTok{!}\FunctionTok{is.na}\NormalTok{(continent))}

\CommentTok{\# 4. Créer le treemap}
\FunctionTok{ggplot}\NormalTok{(usa\_exports\_1900, }\FunctionTok{aes}\NormalTok{(}
  \AttributeTok{area =}\NormalTok{ exports, }\AttributeTok{fill =}\NormalTok{ continent, }\AttributeTok{label =}\NormalTok{ country2, }\AttributeTok{subgroup =}\NormalTok{ continent}
\NormalTok{)) }\SpecialCharTok{+}
  \FunctionTok{geom\_treemap}\NormalTok{() }\SpecialCharTok{+}
  \FunctionTok{geom\_treemap\_text}\NormalTok{(}\AttributeTok{colour =} \StringTok{"white"}\NormalTok{, }\AttributeTok{place =} \StringTok{"centre"}\NormalTok{, }\AttributeTok{grow =} \ConstantTok{TRUE}\NormalTok{, }\AttributeTok{reflow =} \ConstantTok{TRUE}\NormalTok{) }\SpecialCharTok{+}
  \FunctionTok{geom\_treemap\_subgroup\_border}\NormalTok{(}\AttributeTok{color =} \StringTok{"black"}\NormalTok{) }\SpecialCharTok{+}
  \FunctionTok{geom\_treemap\_subgroup\_text}\NormalTok{(}\AttributeTok{place =} \StringTok{"centre"}\NormalTok{, }\AttributeTok{grow =} \ConstantTok{TRUE}\NormalTok{, }\AttributeTok{alpha =} \FloatTok{0.5}\NormalTok{, }\AttributeTok{colour =} \StringTok{"black"}\NormalTok{) }\SpecialCharTok{+}
  \FunctionTok{labs}\NormalTok{(}\AttributeTok{title =} \StringTok{"Treemap des exportations américaines en 1900"}\NormalTok{, }\AttributeTok{fill =} \StringTok{"Continent"}\NormalTok{) }\SpecialCharTok{+}
  \FunctionTok{theme}\NormalTok{(}\AttributeTok{legend.position =} \StringTok{"bottom"}\NormalTok{)}
\end{Highlighting}
\end{Shaded}

\pandocbounded{\includegraphics[keepaspectratio]{Examen_visualisation_trade-networks_files/figure-latex/unnamed-chunk-4-1.pdf}}

\begin{Shaded}
\begin{Highlighting}[]
\CommentTok{\# 1. Sélection des 20 principaux pays importateurs par valeur exportée}
\NormalTok{top20\_usa\_exports }\OtherTok{\textless{}{-}}\NormalTok{ usa\_exports\_1900 }\SpecialCharTok{\%\textgreater{}\%}
  \FunctionTok{arrange}\NormalTok{(}\FunctionTok{desc}\NormalTok{(exports)) }\SpecialCharTok{\%\textgreater{}\%}
  \FunctionTok{slice}\NormalTok{(}\DecValTok{1}\SpecialCharTok{:}\DecValTok{20}\NormalTok{)}

\CommentTok{\# Exporter le tableau en Excel}
\FunctionTok{write.xlsx}\NormalTok{(top20\_usa\_exports, }\StringTok{"data/top20\_usa\_exports.xlsx"}\NormalTok{)}

\CommentTok{\# 2. Treemap avec continents et pays (Top 20)}
\FunctionTok{ggplot}\NormalTok{(top20\_usa\_exports, }\FunctionTok{aes}\NormalTok{(}
  \AttributeTok{area =}\NormalTok{ exports, }\AttributeTok{fill =}\NormalTok{ continent, }\AttributeTok{label =}\NormalTok{ country2, }\AttributeTok{subgroup =}\NormalTok{ continent}
\NormalTok{)) }\SpecialCharTok{+}
  \FunctionTok{geom\_treemap}\NormalTok{() }\SpecialCharTok{+}
  \FunctionTok{geom\_treemap\_text}\NormalTok{(}\AttributeTok{colour =} \StringTok{"white"}\NormalTok{, }\AttributeTok{place =} \StringTok{"centre"}\NormalTok{, }\AttributeTok{grow =} \ConstantTok{TRUE}\NormalTok{, }\AttributeTok{reflow =} \ConstantTok{TRUE}\NormalTok{) }\SpecialCharTok{+}
  \FunctionTok{geom\_treemap\_subgroup\_border}\NormalTok{(}\AttributeTok{color =} \StringTok{"black"}\NormalTok{) }\SpecialCharTok{+}
  \FunctionTok{geom\_treemap\_subgroup\_text}\NormalTok{(}\AttributeTok{place =} \StringTok{"centre"}\NormalTok{, }\AttributeTok{grow =} \ConstantTok{TRUE}\NormalTok{, }\AttributeTok{alpha =} \FloatTok{0.5}\NormalTok{, }\AttributeTok{colour =} \StringTok{"black"}\NormalTok{) }\SpecialCharTok{+}
  \FunctionTok{labs}\NormalTok{(}\AttributeTok{title =} \StringTok{"Top 20 des pays importateurs des États{-}Unis en 1900 (Treemap)"}\NormalTok{,}
       \AttributeTok{fill =} \StringTok{"Continent"}\NormalTok{) }\SpecialCharTok{+}
  \FunctionTok{theme}\NormalTok{(}\AttributeTok{legend.position =} \StringTok{"bottom"}\NormalTok{)}
\end{Highlighting}
\end{Shaded}

\pandocbounded{\includegraphics[keepaspectratio]{Examen_visualisation_trade-networks_files/figure-latex/unnamed-chunk-5-1.pdf}}

\begin{Shaded}
\begin{Highlighting}[]
\FunctionTok{library}\NormalTok{(openxlsx)}
\end{Highlighting}
\end{Shaded}

\subsection{Question 5}\label{question-5}

Maintenant, considérer les données des années 1900 et 2009. Comparer à
partir d'un treemap l'évolution du commerce international entre ces deux
dates. Utiliser ecore une fois l'indicateurs du pourcentage des
exportations. Faites attention, entre temps certains pays int disparu et
d;autres sont apparus. Il faut les inclure correctement dans leur
continent.

\begin{Shaded}
\begin{Highlighting}[]
\CommentTok{\# Charger les bibliothèques utiles}
\FunctionTok{library}\NormalTok{(dplyr)}
\FunctionTok{library}\NormalTok{(ggplot2)}
\FunctionTok{library}\NormalTok{(treemapify)}
\FunctionTok{library}\NormalTok{(countrycode)}

\CommentTok{\# 1. Filtrer les années 1900 et 2009}
\NormalTok{data\_2years }\OtherTok{\textless{}{-}}\NormalTok{ trade }\SpecialCharTok{\%\textgreater{}\%} \FunctionTok{filter}\NormalTok{(year }\SpecialCharTok{\%in\%} \FunctionTok{c}\NormalTok{(}\DecValTok{1900}\NormalTok{, }\DecValTok{2009}\NormalTok{))}

\CommentTok{\# 2. Remplacer les NA par 0 pour éviter les erreurs}
\NormalTok{data\_2years}\SpecialCharTok{$}\NormalTok{exports[}\FunctionTok{is.na}\NormalTok{(data\_2years}\SpecialCharTok{$}\NormalTok{exports)] }\OtherTok{\textless{}{-}} \DecValTok{0}

\CommentTok{\# 3. Calculer les exportations totales par pays exportateur pour chaque année}
\NormalTok{export\_by\_country }\OtherTok{\textless{}{-}}\NormalTok{ data\_2years }\SpecialCharTok{\%\textgreater{}\%}
  \FunctionTok{group\_by}\NormalTok{(year, country1) }\SpecialCharTok{\%\textgreater{}\%}
  \FunctionTok{summarise}\NormalTok{(}\AttributeTok{total\_exports =} \FunctionTok{sum}\NormalTok{(exports), }\AttributeTok{.groups =} \StringTok{"drop"}\NormalTok{)}

\CommentTok{\# 4. Ajouter le pourcentage d\textquotesingle{}exportation mondiale pour chaque année}
\NormalTok{export\_by\_country }\OtherTok{\textless{}{-}}\NormalTok{ export\_by\_country }\SpecialCharTok{\%\textgreater{}\%}
  \FunctionTok{group\_by}\NormalTok{(year) }\SpecialCharTok{\%\textgreater{}\%}
  \FunctionTok{mutate}\NormalTok{(}\AttributeTok{share =}\NormalTok{ total\_exports }\SpecialCharTok{/} \FunctionTok{sum}\NormalTok{(total\_exports))}

\CommentTok{\# 5. Ajouter le continent correspondant à chaque pays exportateur}
\NormalTok{export\_by\_country }\OtherTok{\textless{}{-}}\NormalTok{ export\_by\_country }\SpecialCharTok{\%\textgreater{}\%}
  \FunctionTok{mutate}\NormalTok{(}
    \AttributeTok{continent =} \FunctionTok{countrycode}\NormalTok{(country1, }\AttributeTok{origin =} \StringTok{"country.name"}\NormalTok{, }\AttributeTok{destination =} \StringTok{"continent"}\NormalTok{),}
    \AttributeTok{country1 =} \FunctionTok{case\_when}\NormalTok{(}
\NormalTok{      country1 }\SpecialCharTok{==} \StringTok{"Austria{-}Hungary"} \SpecialCharTok{\textasciitilde{}} \StringTok{"Austria"}\NormalTok{,}
\NormalTok{      country1 }\SpecialCharTok{==} \StringTok{"Yugoslavia"} \SpecialCharTok{\textasciitilde{}} \StringTok{"Serbia"}\NormalTok{,}
      \ConstantTok{TRUE} \SpecialCharTok{\textasciitilde{}}\NormalTok{ country1}
\NormalTok{    )}
\NormalTok{  )}
\end{Highlighting}
\end{Shaded}

\begin{verbatim}
## Warning: There were 2 warnings in `mutate()`.
## The first warning was:
## i In argument: `continent = countrycode(country1, origin = "country.name",
##   destination = "continent")`.
## i In group 1: `year = 1900`.
## Caused by warning:
## ! Some values were not matched unambiguously: Austria-Hungary, Yugoslavia
## i Run `dplyr::last_dplyr_warnings()` to see the 1 remaining warning.
\end{verbatim}

\begin{Shaded}
\begin{Highlighting}[]
\CommentTok{\# Supprimer les lignes avec continent inconnu}
\NormalTok{export\_by\_country }\OtherTok{\textless{}{-}}\NormalTok{ export\_by\_country }\SpecialCharTok{\%\textgreater{}\%} \FunctionTok{filter}\NormalTok{(}\SpecialCharTok{!}\FunctionTok{is.na}\NormalTok{(continent))}

\CommentTok{\# 6. Créer le treemap comparatif}
\FunctionTok{ggplot}\NormalTok{(export\_by\_country, }\FunctionTok{aes}\NormalTok{(}
  \AttributeTok{area =}\NormalTok{ share, }\AttributeTok{fill =}\NormalTok{ continent, }\AttributeTok{label =}\NormalTok{ country1, }\AttributeTok{subgroup =}\NormalTok{ continent}
\NormalTok{)) }\SpecialCharTok{+}
  \FunctionTok{geom\_treemap}\NormalTok{() }\SpecialCharTok{+}
  \FunctionTok{geom\_treemap\_text}\NormalTok{(}\AttributeTok{colour =} \StringTok{"white"}\NormalTok{, }\AttributeTok{place =} \StringTok{"centre"}\NormalTok{, }\AttributeTok{grow =} \ConstantTok{TRUE}\NormalTok{, }\AttributeTok{reflow =} \ConstantTok{TRUE}\NormalTok{) }\SpecialCharTok{+}
  \FunctionTok{geom\_treemap\_subgroup\_border}\NormalTok{(}\AttributeTok{color =} \StringTok{"black"}\NormalTok{) }\SpecialCharTok{+}
  \FunctionTok{geom\_treemap\_subgroup\_text}\NormalTok{(}\AttributeTok{place =} \StringTok{"centre"}\NormalTok{, }\AttributeTok{grow =} \ConstantTok{TRUE}\NormalTok{, }\AttributeTok{alpha =} \FloatTok{0.5}\NormalTok{, }\AttributeTok{colour =} \StringTok{"black"}\NormalTok{) }\SpecialCharTok{+}
  \FunctionTok{facet\_wrap}\NormalTok{(}\SpecialCharTok{\textasciitilde{}}\NormalTok{ year) }\SpecialCharTok{+}
  \FunctionTok{labs}\NormalTok{(}
    \AttributeTok{title =} \StringTok{"Comparaison du commerce international par pays (1900 vs 2009)"}\NormalTok{,}
    \AttributeTok{fill =} \StringTok{"Continent"}
\NormalTok{  ) }\SpecialCharTok{+}
  \FunctionTok{theme}\NormalTok{(}\AttributeTok{legend.position =} \StringTok{"bottom"}\NormalTok{)}
\end{Highlighting}
\end{Shaded}

\pandocbounded{\includegraphics[keepaspectratio]{Examen_visualisation_trade-networks_files/figure-latex/unnamed-chunk-6-1.pdf}}
\#\#\# Interprétation synthétique du treemap comparatif (1900 et 2009)

Le treemap comparatif révèle une transformation profonde du commerce
mondial entre 1900 et 2009. En 1900, les exportations sont fortement
concentrées autour des grandes puissances occidentales (Royaume-Uni,
Allemagne, France, États-Unis), avec quelques partenaires secondaires en
Amérique latine. En 2009, le paysage commercial est nettement plus
multipolaire : aux côtés des États-Unis, on observe la montée en
puissance de la Chine, de l'Allemagne réunifiée, du Japon, de la Corée
du Sud et de plusieurs pays émergents d'Asie et d'Europe de l'Est. Cette
évolution reflète une diversification géographique des échanges, une
réallocation des chaînes de production vers l'Asie, et les conséquences
géopolitiques majeures du XXe siècle (disparition d'États,
indépendances, réunifications). Certains anciens pôles (Portugal,
Belgique, Espagne) reculent en importance relative. En somme, le treemap
met en lumière le passage d'un commerce dominé par quelques empires à un
réseau mondialisé, dense et multipolaire, révélateur de l'essor des
économies émergentes et de la reconfiguration des flux mondiaux.

\subsection{Question 6: Bonus}\label{question-6-bonus}

Les données dont nous disposons sont des données de réseau. Présenter
les graphiques de la page 1433 de l'article. Pour ce faire, suivez les
étapes suivantes:

\begin{enumerate}
\def\labelenumi{\arabic{enumi}.}
\tightlist
\item
  Nous commençons par analyser le commerce international comme un réseau
  non pondéré et dirigé. Pour chaque année de l'ensemble de données,
  créez une matrice d'adjacence dont l'entrée \((i,j)\) est égale à 1 si
  le pays \(i\) exporte vers le pays \(j\). Si cette exportation est
  nulle, alors l'entrée est égale à 0. Nous supposons que les données
  manquantes, indiquées par «NA», représentent un commerce nul. Tracez
  la «densité du réseau», qui est définie au fil du temps comme suit:
\end{enumerate}

\[
    \text{network density}  =  \frac{\text{number of edges}}{\text{number of potential edges}}
  \]

La fonction \texttt{graph.density} du package igraph
(\url{https://igraph.org/}) peut calculer cette mesure à partir d'une
matrice d'adjacence. cet exercice est une exploration dans les données
de réseau avec igraph. Interprétez le résultat.

\begin{enumerate}
\def\labelenumi{\arabic{enumi}.}
\setcounter{enumi}{1}
\item
  Présenter le graphique
\item
  Pour les années 1900, 1955 et 2009, calculez les mesures de centralité
  en fonction du degré, de l'intermédiarité et de la proximité (en
  fonction du degré total) pour chaque année. Pour chaque année, dressez
  la liste des cinq pays qui ont les valeurs les plus élevées de ces
  mesures de centralité. Comment les pays figurant sur les listes
  évoluent-ils au fil du temps? Commentez brièvement les résultats.
\end{enumerate}

\begin{Shaded}
\begin{Highlighting}[]
\DocumentationTok{\#\#\# Partie 1 (densité du réseau)}
\CommentTok{\# Charger les bibliothèques}
\FunctionTok{library}\NormalTok{(igraph)}
\end{Highlighting}
\end{Shaded}

\begin{verbatim}
## 
## Attachement du package : 'igraph'
\end{verbatim}

\begin{verbatim}
## Les objets suivants sont masqués depuis 'package:dplyr':
## 
##     as_data_frame, groups, union
\end{verbatim}

\begin{verbatim}
## Les objets suivants sont masqués depuis 'package:stats':
## 
##     decompose, spectrum
\end{verbatim}

\begin{verbatim}
## L'objet suivant est masqué depuis 'package:base':
## 
##     union
\end{verbatim}

\begin{Shaded}
\begin{Highlighting}[]
\FunctionTok{library}\NormalTok{(dplyr)}
\FunctionTok{library}\NormalTok{(tidyr)}
\end{Highlighting}
\end{Shaded}

\begin{verbatim}
## 
## Attachement du package : 'tidyr'
\end{verbatim}

\begin{verbatim}
## L'objet suivant est masqué depuis 'package:igraph':
## 
##     crossing
\end{verbatim}

\begin{Shaded}
\begin{Highlighting}[]
\FunctionTok{library}\NormalTok{(ggplot2)}


\NormalTok{trade}\SpecialCharTok{$}\NormalTok{exports[}\FunctionTok{is.na}\NormalTok{(trade}\SpecialCharTok{$}\NormalTok{exports)] }\OtherTok{\textless{}{-}} \DecValTok{0}  \CommentTok{\# remplacer les NA par 0}

\CommentTok{\# Liste des années disponibles}
\NormalTok{years }\OtherTok{\textless{}{-}} \FunctionTok{sort}\NormalTok{(}\FunctionTok{unique}\NormalTok{(trade}\SpecialCharTok{$}\NormalTok{year))}

\CommentTok{\# Initialiser un tableau pour stocker les densités}
\NormalTok{densities }\OtherTok{\textless{}{-}} \FunctionTok{data.frame}\NormalTok{()}

\ControlFlowTok{for}\NormalTok{ (yr }\ControlFlowTok{in}\NormalTok{ years) \{}
  \CommentTok{\# Filtrer les données pour l\textquotesingle{}année en cours}
\NormalTok{  data\_year }\OtherTok{\textless{}{-}}\NormalTok{ trade }\SpecialCharTok{\%\textgreater{}\%} \FunctionTok{filter}\NormalTok{(year }\SpecialCharTok{==}\NormalTok{ yr)}
  
  \CommentTok{\# Créer la matrice d\textquotesingle{}adjacence binaire : 1 si export \textgreater{} 0}
\NormalTok{  edge\_list }\OtherTok{\textless{}{-}}\NormalTok{ data\_year }\SpecialCharTok{\%\textgreater{}\%}
    \FunctionTok{filter}\NormalTok{(exports }\SpecialCharTok{\textgreater{}} \DecValTok{0}\NormalTok{) }\SpecialCharTok{\%\textgreater{}\%}
    \FunctionTok{select}\NormalTok{(}\AttributeTok{from =}\NormalTok{ country1, }\AttributeTok{to =}\NormalTok{ country2)}
  
  \CommentTok{\# Créer le graphe dirigé}
\NormalTok{  g }\OtherTok{\textless{}{-}} \FunctionTok{graph\_from\_data\_frame}\NormalTok{(edge\_list, }\AttributeTok{directed =} \ConstantTok{TRUE}\NormalTok{)}
  
  \CommentTok{\# Calculer la densité du réseau}
\NormalTok{  dens }\OtherTok{\textless{}{-}} \FunctionTok{edge\_density}\NormalTok{(g)}
  
  \CommentTok{\# Ajouter au tableau}
\NormalTok{  densities }\OtherTok{\textless{}{-}} \FunctionTok{rbind}\NormalTok{(densities, }\FunctionTok{data.frame}\NormalTok{(}\AttributeTok{year =}\NormalTok{ yr, }\AttributeTok{density =}\NormalTok{ dens))}
\NormalTok{\}}

\CommentTok{\# Visualiser la densité du réseau au fil du temps}
\FunctionTok{ggplot}\NormalTok{(densities, }\FunctionTok{aes}\NormalTok{(}\AttributeTok{x =}\NormalTok{ year, }\AttributeTok{y =}\NormalTok{ density)) }\SpecialCharTok{+}
  \FunctionTok{geom\_line}\NormalTok{(}\AttributeTok{color =} \StringTok{"steelblue"}\NormalTok{, }\AttributeTok{size =} \FloatTok{1.2}\NormalTok{) }\SpecialCharTok{+}
  \FunctionTok{geom\_point}\NormalTok{(}\AttributeTok{color =} \StringTok{"darkred"}\NormalTok{, }\AttributeTok{size =} \DecValTok{2}\NormalTok{) }\SpecialCharTok{+}
  \FunctionTok{labs}\NormalTok{(}\AttributeTok{title =} \StringTok{"Densité du réseau commercial international (1900–2009)"}\NormalTok{,}
       \AttributeTok{x =} \StringTok{"Année"}\NormalTok{, }\AttributeTok{y =} \StringTok{"Densité du réseau"}\NormalTok{) }\SpecialCharTok{+}
  \FunctionTok{theme\_minimal}\NormalTok{()}
\end{Highlighting}
\end{Shaded}

\begin{verbatim}
## Warning: Using `size` aesthetic for lines was deprecated in ggplot2 3.4.0.
## i Please use `linewidth` instead.
## This warning is displayed once every 8 hours.
## Call `lifecycle::last_lifecycle_warnings()` to see where this warning was
## generated.
\end{verbatim}

\pandocbounded{\includegraphics[keepaspectratio]{Examen_visualisation_trade-networks_files/figure-latex/unnamed-chunk-7-1.pdf}}

\begin{Shaded}
\begin{Highlighting}[]
\DocumentationTok{\#\#\# Autre graphique}

\FunctionTok{library}\NormalTok{(ggplot2)}

\FunctionTok{ggplot}\NormalTok{(densities, }\FunctionTok{aes}\NormalTok{(}\AttributeTok{x =}\NormalTok{ year, }\AttributeTok{y =}\NormalTok{ density)) }\SpecialCharTok{+}
  \FunctionTok{geom\_line}\NormalTok{(}\AttributeTok{color =} \StringTok{"steelblue"}\NormalTok{, }\AttributeTok{size =} \FloatTok{1.3}\NormalTok{) }\SpecialCharTok{+}
  \FunctionTok{geom\_point}\NormalTok{(}\AttributeTok{color =} \StringTok{"darkred"}\NormalTok{, }\AttributeTok{size =} \DecValTok{3}\NormalTok{) }\SpecialCharTok{+}
  \FunctionTok{geom\_text}\NormalTok{(}\FunctionTok{aes}\NormalTok{(}\AttributeTok{label =} \FunctionTok{round}\NormalTok{(density, }\DecValTok{3}\NormalTok{)), }\AttributeTok{vjust =} \SpecialCharTok{{-}}\FloatTok{0.8}\NormalTok{, }\AttributeTok{size =} \DecValTok{4}\NormalTok{) }\SpecialCharTok{+}
  \FunctionTok{labs}\NormalTok{(}\AttributeTok{title =} \StringTok{"Évolution de la densité du réseau commercial mondial"}\NormalTok{,}
       \AttributeTok{x =} \StringTok{"Année"}\NormalTok{, }\AttributeTok{y =} \StringTok{"Densité du réseau"}\NormalTok{) }\SpecialCharTok{+}
  \FunctionTok{theme\_minimal}\NormalTok{(}\AttributeTok{base\_size =} \DecValTok{14}\NormalTok{)}
\end{Highlighting}
\end{Shaded}

\pandocbounded{\includegraphics[keepaspectratio]{Examen_visualisation_trade-networks_files/figure-latex/unnamed-chunk-8-1.pdf}}

\begin{Shaded}
\begin{Highlighting}[]
\DocumentationTok{\#\#\# Barplot clair par année}

\FunctionTok{ggplot}\NormalTok{(densities, }\FunctionTok{aes}\NormalTok{(}\AttributeTok{x =} \FunctionTok{factor}\NormalTok{(year), }\AttributeTok{y =}\NormalTok{ density, }\AttributeTok{fill =}\NormalTok{ density)) }\SpecialCharTok{+}
  \FunctionTok{geom\_col}\NormalTok{(}\AttributeTok{show.legend =} \ConstantTok{FALSE}\NormalTok{) }\SpecialCharTok{+}
  \FunctionTok{geom\_text}\NormalTok{(}\FunctionTok{aes}\NormalTok{(}\AttributeTok{label =} \FunctionTok{round}\NormalTok{(density, }\DecValTok{3}\NormalTok{)), }\AttributeTok{vjust =} \SpecialCharTok{{-}}\FloatTok{0.5}\NormalTok{, }\AttributeTok{size =} \DecValTok{4}\NormalTok{) }\SpecialCharTok{+}
  \FunctionTok{scale\_fill\_gradient}\NormalTok{(}\AttributeTok{low =} \StringTok{"\#FFDDC1"}\NormalTok{, }\AttributeTok{high =} \StringTok{"\#2C3E50"}\NormalTok{) }\SpecialCharTok{+}
  \FunctionTok{labs}\NormalTok{(}\AttributeTok{title =} \StringTok{"Densité du réseau commercial international (par année)"}\NormalTok{,}
       \AttributeTok{x =} \StringTok{"Année"}\NormalTok{, }\AttributeTok{y =} \StringTok{"Densité"}\NormalTok{) }\SpecialCharTok{+}
  \FunctionTok{theme\_minimal}\NormalTok{(}\AttributeTok{base\_size =} \DecValTok{14}\NormalTok{)}
\end{Highlighting}
\end{Shaded}

\pandocbounded{\includegraphics[keepaspectratio]{Examen_visualisation_trade-networks_files/figure-latex/unnamed-chunk-9-1.pdf}}

\begin{Shaded}
\begin{Highlighting}[]
\DocumentationTok{\#\#\# Mesures de centralité réseau}
\CommentTok{\# Années à analyser}
\NormalTok{target\_years }\OtherTok{\textless{}{-}} \FunctionTok{c}\NormalTok{(}\DecValTok{1900}\NormalTok{, }\DecValTok{1955}\NormalTok{, }\DecValTok{2009}\NormalTok{)}

\CommentTok{\# Initialiser un tableau de résultats}
\NormalTok{centrality\_results }\OtherTok{\textless{}{-}} \FunctionTok{list}\NormalTok{()}

\CommentTok{\# Boucle sur les années cibles}
\ControlFlowTok{for}\NormalTok{ (yr }\ControlFlowTok{in}\NormalTok{ target\_years) \{}
  \CommentTok{\# Filtrer les données pour l’année}
\NormalTok{  data\_year }\OtherTok{\textless{}{-}}\NormalTok{ trade }\SpecialCharTok{\%\textgreater{}\%}
    \FunctionTok{filter}\NormalTok{(year }\SpecialCharTok{==}\NormalTok{ yr }\SpecialCharTok{\&}\NormalTok{ exports }\SpecialCharTok{\textgreater{}} \DecValTok{0}\NormalTok{) }\SpecialCharTok{\%\textgreater{}\%}
    \FunctionTok{select}\NormalTok{(}\AttributeTok{from =}\NormalTok{ country1, }\AttributeTok{to =}\NormalTok{ country2)}
  
  \CommentTok{\# Construire le graphe dirigé}
\NormalTok{  g }\OtherTok{\textless{}{-}} \FunctionTok{graph\_from\_data\_frame}\NormalTok{(data\_year, }\AttributeTok{directed =} \ConstantTok{TRUE}\NormalTok{)}
  
  \CommentTok{\# Calculer les centralités}
\NormalTok{  degree\_cent }\OtherTok{\textless{}{-}} \FunctionTok{degree}\NormalTok{(g, }\AttributeTok{mode =} \StringTok{"all"}\NormalTok{)}
\NormalTok{  between\_cent }\OtherTok{\textless{}{-}} \FunctionTok{betweenness}\NormalTok{(g, }\AttributeTok{normalized =} \ConstantTok{TRUE}\NormalTok{)}
\NormalTok{  close\_cent }\OtherTok{\textless{}{-}} \FunctionTok{closeness}\NormalTok{(g, }\AttributeTok{normalized =} \ConstantTok{TRUE}\NormalTok{)}
  
  \CommentTok{\# Créer une table avec les résultats}
\NormalTok{  centrality\_df }\OtherTok{\textless{}{-}} \FunctionTok{data.frame}\NormalTok{(}
    \AttributeTok{country =} \FunctionTok{names}\NormalTok{(degree\_cent),}
    \AttributeTok{year =}\NormalTok{ yr,}
    \AttributeTok{degree =}\NormalTok{ degree\_cent,}
    \AttributeTok{betweenness =}\NormalTok{ between\_cent,}
    \AttributeTok{closeness =}\NormalTok{ close\_cent}
\NormalTok{  )}
  
  \CommentTok{\# Garder seulement les top 5 pour chaque mesure}
\NormalTok{  top5\_degree }\OtherTok{\textless{}{-}}\NormalTok{ centrality\_df }\SpecialCharTok{\%\textgreater{}\%} \FunctionTok{arrange}\NormalTok{(}\FunctionTok{desc}\NormalTok{(degree)) }\SpecialCharTok{\%\textgreater{}\%} \FunctionTok{slice}\NormalTok{(}\DecValTok{1}\SpecialCharTok{:}\DecValTok{5}\NormalTok{)}
\NormalTok{  top5\_between }\OtherTok{\textless{}{-}}\NormalTok{ centrality\_df }\SpecialCharTok{\%\textgreater{}\%} \FunctionTok{arrange}\NormalTok{(}\FunctionTok{desc}\NormalTok{(betweenness)) }\SpecialCharTok{\%\textgreater{}\%} \FunctionTok{slice}\NormalTok{(}\DecValTok{1}\SpecialCharTok{:}\DecValTok{5}\NormalTok{)}
\NormalTok{  top5\_close }\OtherTok{\textless{}{-}}\NormalTok{ centrality\_df }\SpecialCharTok{\%\textgreater{}\%} \FunctionTok{arrange}\NormalTok{(}\FunctionTok{desc}\NormalTok{(closeness)) }\SpecialCharTok{\%\textgreater{}\%} \FunctionTok{slice}\NormalTok{(}\DecValTok{1}\SpecialCharTok{:}\DecValTok{5}\NormalTok{)}
  
  \CommentTok{\# Stocker dans une liste}
\NormalTok{  centrality\_results[[}\FunctionTok{as.character}\NormalTok{(yr)]] }\OtherTok{\textless{}{-}} \FunctionTok{list}\NormalTok{(}
    \AttributeTok{degree =}\NormalTok{ top5\_degree,}
    \AttributeTok{betweenness =}\NormalTok{ top5\_between,}
    \AttributeTok{closeness =}\NormalTok{ top5\_close}
\NormalTok{  )}
\NormalTok{\}}

\CommentTok{\# Affichage des résultats}
\NormalTok{centrality\_results}
\end{Highlighting}
\end{Shaded}

\begin{verbatim}
## $`1900`
## $`1900`$degree
##                                           country year degree betweenness
## United Kingdom                     United Kingdom 1900     65  0.35248336
## United States of America United States of America 1900     54  0.11185210
## France                                     France 1900     47  0.05046188
## Germany                                   Germany 1900     47  0.05360914
## Belgium                                   Belgium 1900     42  0.03145395
##                          closeness
## United Kingdom           0.9459459
## United States of America 0.8333333
## France                   0.7608696
## Germany                  0.7777778
## Belgium                  0.7142857
## 
## $`1900`$betweenness
##                                           country year degree betweenness
## United Kingdom                     United Kingdom 1900     65  0.35248336
## United States of America United States of America 1900     54  0.11185210
## Japan                                       Japan 1900     28  0.07160401
## Austria-Hungary                   Austria-Hungary 1900     33  0.06547755
## Germany                                   Germany 1900     47  0.05360914
##                          closeness
## United Kingdom           0.9459459
## United States of America 0.8333333
## Japan                    0.6250000
## Austria-Hungary          0.6603774
## Germany                  0.7777778
## 
## $`1900`$closeness
##                                           country year degree betweenness
## United Kingdom                     United Kingdom 1900     65  0.35248336
## United States of America United States of America 1900     54  0.11185210
## Germany                                   Germany 1900     47  0.05360914
## France                                     France 1900     47  0.05046188
## Belgium                                   Belgium 1900     42  0.03145395
##                          closeness
## United Kingdom           0.9459459
## United States of America 0.8333333
## Germany                  0.7777778
## France                   0.7608696
## Belgium                  0.7142857
## 
## 
## $`1955`
## $`1955`$degree
##                                         country year degree betweenness
## United Kingdom                   United Kingdom 1955    151  0.04818518
## German Federal Republic German Federal Republic 1955    147  0.03415625
## Netherlands                         Netherlands 1955    146  0.02247210
## Italy                                     Italy 1955    146  0.03972992
## France                                   France 1955    145  0.02497874
##                         closeness
## United Kingdom          0.9634146
## German Federal Republic 0.9518072
## Netherlands             0.9634146
## Italy                   0.9634146
## France                  0.9518072
## 
## $`1955`$betweenness
##                                           country year degree betweenness
## United Kingdom                     United Kingdom 1955    151  0.04818518
## Italy                                       Italy 1955    146  0.03972992
## United States of America United States of America 1955    142  0.03583847
## German Federal Republic   German Federal Republic 1955    147  0.03415625
## France                                     France 1955    145  0.02497874
##                          closeness
## United Kingdom           0.9634146
## Italy                    0.9634146
## United States of America 0.9404762
## German Federal Republic  0.9518072
## France                   0.9518072
## 
## $`1955`$closeness
##                                         country year degree betweenness
## United Kingdom                   United Kingdom 1955    151  0.04818518
## Netherlands                         Netherlands 1955    146  0.02247210
## Italy                                     Italy 1955    146  0.03972992
## France                                   France 1955    145  0.02497874
## German Federal Republic German Federal Republic 1955    147  0.03415625
##                         closeness
## United Kingdom          0.9634146
## Netherlands             0.9634146
## Italy                   0.9634146
## France                  0.9518072
## German Federal Republic 0.9518072
## 
## 
## $`2009`
## $`2009`$degree
##                                           country year degree betweenness
## United Kingdom                     United Kingdom 2009    363  0.01266067
## China                                       China 2009    363  0.01325424
## India                                       India 2009    362  0.01330262
## United States of America United States of America 2009    361  0.01498608
## France                                     France 2009    361  0.01191885
##                          closeness
## United Kingdom           0.9739583
## China                    0.9790576
## India                    0.9790576
## United States of America 0.9739583
## France                   0.9739583
## 
## $`2009`$betweenness
##                                           country year degree betweenness
## Taiwan                                     Taiwan 2009     42  0.02133333
## United States of America United States of America 2009    361  0.01498608
## Canada                                     Canada 2009    358  0.01396141
## Japan                                       Japan 2009    361  0.01361387
## India                                       India 2009    362  0.01330262
##                          closeness
## Taiwan                   0.5297450
## United States of America 0.9739583
## Canada                   0.9639175
## Japan                    0.9739583
## India                    0.9790576
## 
## $`2009`$closeness
##                                           country year degree betweenness
## China                                       China 2009    363  0.01325424
## India                                       India 2009    362  0.01330262
## United States of America United States of America 2009    361  0.01498608
## United Kingdom                     United Kingdom 2009    363  0.01266067
## France                                     France 2009    361  0.01191885
##                          closeness
## China                    0.9790576
## India                    0.9790576
## United States of America 0.9739583
## United Kingdom           0.9739583
## France                   0.9739583
\end{verbatim}

\begin{Shaded}
\begin{Highlighting}[]
\FunctionTok{library}\NormalTok{(tidyverse)}
\end{Highlighting}
\end{Shaded}

\begin{verbatim}
## -- Attaching core tidyverse packages ------------------------ tidyverse 2.0.0 --
## v forcats   1.0.0     v stringr   1.5.1
## v lubridate 1.9.4     v tibble    3.3.0
## v purrr     1.1.0     
## -- Conflicts ------------------------------------------ tidyverse_conflicts() --
## x lubridate::%--%()       masks igraph::%--%()
## x tibble::as_data_frame() masks igraph::as_data_frame(), dplyr::as_data_frame()
## x purrr::compose()        masks igraph::compose()
## x tidyr::crossing()       masks igraph::crossing()
## x dplyr::filter()         masks stats::filter()
## x dplyr::lag()            masks stats::lag()
## x purrr::simplify()       masks igraph::simplify()
## i Use the conflicted package (<http://conflicted.r-lib.org/>) to force all conflicts to become errors
\end{verbatim}

\begin{Shaded}
\begin{Highlighting}[]
\FunctionTok{library}\NormalTok{(igraph)}

\CommentTok{\# 🧹 Étape 1 : Nettoyage des données (remplace NA par 0 dans exports)}
\NormalTok{trade\_clean }\OtherTok{\textless{}{-}}\NormalTok{ trade }\SpecialCharTok{\%\textgreater{}\%}
  \FunctionTok{select}\NormalTok{(country1, country2, year, exports) }\SpecialCharTok{\%\textgreater{}\%}
  \FunctionTok{mutate}\NormalTok{(}\AttributeTok{exports =} \FunctionTok{ifelse}\NormalTok{(}\FunctionTok{is.na}\NormalTok{(exports), }\DecValTok{0}\NormalTok{, exports))}
\FunctionTok{library}\NormalTok{(tidyverse)}
\FunctionTok{library}\NormalTok{(igraph)}

\CommentTok{\# 🧠 Étape 2 corrigée : calcul densité avec gestion des doublons}
\NormalTok{densites }\OtherTok{\textless{}{-}}\NormalTok{ trade\_clean }\SpecialCharTok{\%\textgreater{}\%}
  \FunctionTok{group\_by}\NormalTok{(year) }\SpecialCharTok{\%\textgreater{}\%}
  \FunctionTok{group\_split}\NormalTok{() }\SpecialCharTok{\%\textgreater{}\%}
  \FunctionTok{map\_dfr}\NormalTok{(}\ControlFlowTok{function}\NormalTok{(df\_year) \{}
\NormalTok{    annee }\OtherTok{\textless{}{-}} \FunctionTok{unique}\NormalTok{(df\_year}\SpecialCharTok{$}\NormalTok{year)}
    
    \CommentTok{\# 1. Créer valeur binaire}
\NormalTok{    df\_bin }\OtherTok{\textless{}{-}}\NormalTok{ df\_year }\SpecialCharTok{\%\textgreater{}\%}
      \FunctionTok{mutate}\NormalTok{(}\AttributeTok{value =} \FunctionTok{ifelse}\NormalTok{(exports }\SpecialCharTok{\textgreater{}} \DecValTok{0}\NormalTok{, }\DecValTok{1}\NormalTok{, }\DecValTok{0}\NormalTok{)) }\SpecialCharTok{\%\textgreater{}\%}
      \FunctionTok{group\_by}\NormalTok{(country1, country2) }\SpecialCharTok{\%\textgreater{}\%}         \CommentTok{\# ⚠️ Aggrégation si doublons}
      \FunctionTok{summarise}\NormalTok{(}\AttributeTok{value =} \FunctionTok{max}\NormalTok{(value), }\AttributeTok{.groups =} \StringTok{"drop"}\NormalTok{)}
    
    \CommentTok{\# 2. Création matrice d’adjacence}
\NormalTok{    mat }\OtherTok{\textless{}{-}}\NormalTok{ df\_bin }\SpecialCharTok{\%\textgreater{}\%}
      \FunctionTok{pivot\_wider}\NormalTok{(}\AttributeTok{names\_from =}\NormalTok{ country2, }\AttributeTok{values\_from =}\NormalTok{ value, }\AttributeTok{values\_fill =} \DecValTok{0}\NormalTok{) }\SpecialCharTok{\%\textgreater{}\%}
      \FunctionTok{column\_to\_rownames}\NormalTok{(}\StringTok{"country1"}\NormalTok{) }\SpecialCharTok{\%\textgreater{}\%}
      \FunctionTok{as.matrix}\NormalTok{()}
    
    \CommentTok{\# 3. Graphe dirigé}
\NormalTok{    g }\OtherTok{\textless{}{-}} \FunctionTok{graph\_from\_adjacency\_matrix}\NormalTok{(mat, }\AttributeTok{mode =} \StringTok{"directed"}\NormalTok{, }\AttributeTok{diag =} \ConstantTok{FALSE}\NormalTok{)}
    
    \CommentTok{\# 4. Densité du graphe}
    \FunctionTok{tibble}\NormalTok{(}\AttributeTok{annee =}\NormalTok{ annee, }\AttributeTok{densite =} \FunctionTok{graph.density}\NormalTok{(g))}
\NormalTok{  \})}
\end{Highlighting}
\end{Shaded}

\begin{verbatim}
## Warning: `graph.density()` was deprecated in igraph 2.0.0.
## i Please use `edge_density()` instead.
## This warning is displayed once every 8 hours.
## Call `lifecycle::last_lifecycle_warnings()` to see where this warning was
## generated.
\end{verbatim}

\begin{Shaded}
\begin{Highlighting}[]
\FunctionTok{ggplot}\NormalTok{(densites, }\FunctionTok{aes}\NormalTok{(}\AttributeTok{x =}\NormalTok{ annee, }\AttributeTok{y =}\NormalTok{ densite)) }\SpecialCharTok{+}
  \FunctionTok{geom\_line}\NormalTok{(}\AttributeTok{color =} \StringTok{"\#2b8cbe"}\NormalTok{, }\AttributeTok{size =} \FloatTok{1.2}\NormalTok{) }\SpecialCharTok{+}
  \FunctionTok{geom\_point}\NormalTok{(}\AttributeTok{color =} \StringTok{"\#f03b20"}\NormalTok{, }\AttributeTok{size =} \DecValTok{2}\NormalTok{) }\SpecialCharTok{+}
  \FunctionTok{theme\_minimal}\NormalTok{(}\AttributeTok{base\_size =} \DecValTok{14}\NormalTok{) }\SpecialCharTok{+}
  \FunctionTok{labs}\NormalTok{(}
    \AttributeTok{title =} \StringTok{"📊 Densité du réseau commercial international (1900–2009)"}\NormalTok{,}
    \AttributeTok{x =} \StringTok{"Année"}\NormalTok{,}
    \AttributeTok{y =} \StringTok{"Densité du réseau"}
\NormalTok{  )}
\end{Highlighting}
\end{Shaded}

\pandocbounded{\includegraphics[keepaspectratio]{Examen_visualisation_trade-networks_files/figure-latex/unnamed-chunk-11-1.pdf}}

\begin{Shaded}
\begin{Highlighting}[]
\NormalTok{calculer\_centralites }\OtherTok{\textless{}{-}} \ControlFlowTok{function}\NormalTok{(df, annee\_cible) \{}
  \CommentTok{\# 🔍 Filtrage de l’année cible + binarisation des flux}
\NormalTok{  df\_filtre }\OtherTok{\textless{}{-}}\NormalTok{ df }\SpecialCharTok{\%\textgreater{}\%}
    \FunctionTok{filter}\NormalTok{(year }\SpecialCharTok{==}\NormalTok{ annee\_cible) }\SpecialCharTok{\%\textgreater{}\%}
    \FunctionTok{mutate}\NormalTok{(}\AttributeTok{value =} \FunctionTok{ifelse}\NormalTok{(exports }\SpecialCharTok{\textgreater{}} \DecValTok{0}\NormalTok{, }\DecValTok{1}\NormalTok{, }\DecValTok{0}\NormalTok{)) }\SpecialCharTok{\%\textgreater{}\%}
    \FunctionTok{group\_by}\NormalTok{(country1, country2) }\SpecialCharTok{\%\textgreater{}\%}
    \FunctionTok{summarise}\NormalTok{(}\AttributeTok{value =} \FunctionTok{max}\NormalTok{(value), }\AttributeTok{.groups =} \StringTok{"drop"}\NormalTok{)}
  
  \CommentTok{\# 📊 Création de la matrice d’adjacence}
\NormalTok{  mat }\OtherTok{\textless{}{-}}\NormalTok{ df\_filtre }\SpecialCharTok{\%\textgreater{}\%}
    \FunctionTok{pivot\_wider}\NormalTok{(}\AttributeTok{names\_from =}\NormalTok{ country2, }\AttributeTok{values\_from =}\NormalTok{ value, }\AttributeTok{values\_fill =} \DecValTok{0}\NormalTok{) }\SpecialCharTok{\%\textgreater{}\%}
    \FunctionTok{column\_to\_rownames}\NormalTok{(}\StringTok{"country1"}\NormalTok{) }\SpecialCharTok{\%\textgreater{}\%}
    \FunctionTok{as.matrix}\NormalTok{()}
  
  \CommentTok{\# 🌐 Création du graphe dirigé}
\NormalTok{  g }\OtherTok{\textless{}{-}} \FunctionTok{graph\_from\_adjacency\_matrix}\NormalTok{(mat, }\AttributeTok{mode =} \StringTok{"directed"}\NormalTok{, }\AttributeTok{diag =} \ConstantTok{FALSE}\NormalTok{)}
  
  \CommentTok{\# 📌 Calcul des trois centralités}
\NormalTok{  deg }\OtherTok{\textless{}{-}} \FunctionTok{degree}\NormalTok{(g, }\AttributeTok{mode =} \StringTok{"all"}\NormalTok{)}
\NormalTok{  bet }\OtherTok{\textless{}{-}} \FunctionTok{betweenness}\NormalTok{(g, }\AttributeTok{directed =} \ConstantTok{TRUE}\NormalTok{)}
\NormalTok{  clo }\OtherTok{\textless{}{-}} \FunctionTok{closeness}\NormalTok{(g, }\AttributeTok{mode =} \StringTok{"all"}\NormalTok{)}
  
  \FunctionTok{tibble}\NormalTok{(}\AttributeTok{pays =} \FunctionTok{names}\NormalTok{(deg), }\AttributeTok{degre =}\NormalTok{ deg, }\AttributeTok{intermediarite =}\NormalTok{ bet, }\AttributeTok{proximite =}\NormalTok{ clo) }\SpecialCharTok{\%\textgreater{}\%}
    \FunctionTok{mutate}\NormalTok{(}\AttributeTok{annee =}\NormalTok{ annee\_cible) }\SpecialCharTok{\%\textgreater{}\%}
    \FunctionTok{pivot\_longer}\NormalTok{(}\AttributeTok{cols =} \FunctionTok{c}\NormalTok{(degre, intermediarite, proximite), }\AttributeTok{names\_to =} \StringTok{"mesure"}\NormalTok{, }\AttributeTok{values\_to =} \StringTok{"valeur"}\NormalTok{) }\SpecialCharTok{\%\textgreater{}\%}
    \FunctionTok{group\_by}\NormalTok{(annee, mesure) }\SpecialCharTok{\%\textgreater{}\%}
    \FunctionTok{slice\_max}\NormalTok{(}\AttributeTok{order\_by =}\NormalTok{ valeur, }\AttributeTok{n =} \DecValTok{5}\NormalTok{) }\SpecialCharTok{\%\textgreater{}\%}
    \FunctionTok{arrange}\NormalTok{(annee, mesure, }\FunctionTok{desc}\NormalTok{(valeur))}
\NormalTok{\}}

\NormalTok{resultats\_centralites }\OtherTok{\textless{}{-}} \FunctionTok{map\_dfr}\NormalTok{(}\FunctionTok{c}\NormalTok{(}\DecValTok{1900}\NormalTok{, }\DecValTok{1955}\NormalTok{, }\DecValTok{2009}\NormalTok{), }\SpecialCharTok{\textasciitilde{}}\FunctionTok{calculer\_centralites}\NormalTok{(trade\_clean, .x))}


\CommentTok{\# Exporter le tableau en Excel}
\FunctionTok{write.xlsx}\NormalTok{(resultats\_centralites, }\StringTok{"data/resultats\_centralites.xlsx"}\NormalTok{)}

\NormalTok{resultats\_centralites }\SpecialCharTok{\%\textgreater{}\%}
  \FunctionTok{select}\NormalTok{(annee, mesure, pays, valeur) }\SpecialCharTok{\%\textgreater{}\%}
  \FunctionTok{arrange}\NormalTok{(annee, mesure, }\FunctionTok{desc}\NormalTok{(valeur))}
\end{Highlighting}
\end{Shaded}

\begin{verbatim}
## # A tibble: 46 x 4
## # Groups:   annee, mesure [9]
##    annee mesure         pays                     valeur
##    <dbl> <chr>          <chr>                     <dbl>
##  1  1900 degre          United States of America   57  
##  2  1900 degre          Germany                    44  
##  3  1900 degre          United Kingdom             38  
##  4  1900 degre          Belgium                    36  
##  5  1900 degre          Uruguay                    36  
##  6  1900 intermediarite United States of America  397. 
##  7  1900 intermediarite United Kingdom            168. 
##  8  1900 intermediarite Germany                   103. 
##  9  1900 intermediarite Japan                      75.6
## 10  1900 intermediarite Belgium                    60.1
## # i 36 more rows
\end{verbatim}

\begin{Shaded}
\begin{Highlighting}[]
\FunctionTok{library}\NormalTok{(ggplot2)}

\CommentTok{\# 🧽 Nettoyage pour une présentation claire}
\NormalTok{resultats\_centralites}\SpecialCharTok{$}\NormalTok{mesure }\OtherTok{\textless{}{-}} \FunctionTok{recode}\NormalTok{(resultats\_centralites}\SpecialCharTok{$}\NormalTok{mesure,}
  \AttributeTok{degre =} \StringTok{"🎯 Degré"}\NormalTok{,}
  \AttributeTok{intermediarite =} \StringTok{"🔀 Intermédiarité"}\NormalTok{,}
  \AttributeTok{proximite =} \StringTok{"📍 Proximité"}
\NormalTok{)}

\CommentTok{\# 📊 Graphe en barres, groupé par année et mesure}
\FunctionTok{ggplot}\NormalTok{(resultats\_centralites, }\FunctionTok{aes}\NormalTok{(}\AttributeTok{x =} \FunctionTok{reorder}\NormalTok{(pays, valeur), }\AttributeTok{y =}\NormalTok{ valeur, }\AttributeTok{fill =}\NormalTok{ pays)) }\SpecialCharTok{+}
  \FunctionTok{geom\_col}\NormalTok{(}\AttributeTok{show.legend =} \ConstantTok{FALSE}\NormalTok{) }\SpecialCharTok{+}
  \FunctionTok{facet\_grid}\NormalTok{(mesure }\SpecialCharTok{\textasciitilde{}}\NormalTok{ annee, }\AttributeTok{scales =} \StringTok{"free"}\NormalTok{) }\SpecialCharTok{+}
  \FunctionTok{coord\_flip}\NormalTok{() }\SpecialCharTok{+}
  \FunctionTok{labs}\NormalTok{(}
    \AttributeTok{title =} \StringTok{"🏛️ Mesures de centralité du commerce international (1900, 1955, 2009)"}\NormalTok{,}
    \AttributeTok{x =} \StringTok{"Pays"}\NormalTok{,}
    \AttributeTok{y =} \StringTok{"Valeur de la centralité"}
\NormalTok{  ) }\SpecialCharTok{+}
  \FunctionTok{theme\_minimal}\NormalTok{(}\AttributeTok{base\_size =} \DecValTok{13}\NormalTok{) }\SpecialCharTok{+}
  \FunctionTok{theme}\NormalTok{(}\AttributeTok{strip.text =} \FunctionTok{element\_text}\NormalTok{(}\AttributeTok{face =} \StringTok{"bold"}\NormalTok{, }\AttributeTok{color =} \StringTok{"black"}\NormalTok{))}
\end{Highlighting}
\end{Shaded}

\pandocbounded{\includegraphics[keepaspectratio]{Examen_visualisation_trade-networks_files/figure-latex/unnamed-chunk-13-1.pdf}}

\begin{Shaded}
\begin{Highlighting}[]
\FunctionTok{library}\NormalTok{(igraph)}
\FunctionTok{library}\NormalTok{(tidyverse)}
\FunctionTok{library}\NormalTok{(ggraph)}

\CommentTok{\# 🔁 Fonction pour créer le graphe pour une année spécifique}
\NormalTok{creer\_graphe\_reseau }\OtherTok{\textless{}{-}} \ControlFlowTok{function}\NormalTok{(data, annee) \{}
  \CommentTok{\# Filtrage de l\textquotesingle{}année}
\NormalTok{  df\_annee }\OtherTok{\textless{}{-}}\NormalTok{ data }\SpecialCharTok{\%\textgreater{}\%}
    \FunctionTok{filter}\NormalTok{(year }\SpecialCharTok{==}\NormalTok{ annee) }\SpecialCharTok{\%\textgreater{}\%}
    \FunctionTok{mutate}\NormalTok{(}\AttributeTok{value =} \FunctionTok{ifelse}\NormalTok{(exports }\SpecialCharTok{\textgreater{}} \DecValTok{0}\NormalTok{, }\DecValTok{1}\NormalTok{, }\DecValTok{0}\NormalTok{)) }\SpecialCharTok{\%\textgreater{}\%}
    \FunctionTok{filter}\NormalTok{(value }\SpecialCharTok{==} \DecValTok{1}\NormalTok{)}
  
  \CommentTok{\# Création du graphe dirigé non pondéré}
\NormalTok{  graphe }\OtherTok{\textless{}{-}} \FunctionTok{graph\_from\_data\_frame}\NormalTok{(df\_annee[, }\FunctionTok{c}\NormalTok{(}\StringTok{"country1"}\NormalTok{, }\StringTok{"country2"}\NormalTok{)], }\AttributeTok{directed =} \ConstantTok{TRUE}\NormalTok{)}
  
  \CommentTok{\# Visualisation du graphe avec ggraph}
  \FunctionTok{ggraph}\NormalTok{(graphe, }\AttributeTok{layout =} \StringTok{"fr"}\NormalTok{) }\SpecialCharTok{+}  \CommentTok{\# "fr" = Fruchterman{-}Reingold layout}
    \FunctionTok{geom\_edge\_link}\NormalTok{(}\AttributeTok{alpha =} \FloatTok{0.3}\NormalTok{, }\AttributeTok{arrow =} \FunctionTok{arrow}\NormalTok{(}\AttributeTok{length =} \FunctionTok{unit}\NormalTok{(}\DecValTok{2}\NormalTok{, }\StringTok{\textquotesingle{}mm\textquotesingle{}}\NormalTok{)), }\AttributeTok{end\_cap =} \FunctionTok{circle}\NormalTok{(}\DecValTok{2}\NormalTok{, }\StringTok{\textquotesingle{}mm\textquotesingle{}}\NormalTok{)) }\SpecialCharTok{+}
    \FunctionTok{geom\_node\_point}\NormalTok{(}\AttributeTok{size =} \DecValTok{3}\NormalTok{, }\AttributeTok{color =} \StringTok{"\#2b8cbe"}\NormalTok{) }\SpecialCharTok{+}
    \FunctionTok{geom\_node\_text}\NormalTok{(}\FunctionTok{aes}\NormalTok{(}\AttributeTok{label =}\NormalTok{ name), }\AttributeTok{repel =} \ConstantTok{TRUE}\NormalTok{, }\AttributeTok{size =} \DecValTok{3}\NormalTok{, }\AttributeTok{vjust =} \FloatTok{1.5}\NormalTok{) }\SpecialCharTok{+}
    \FunctionTok{labs}\NormalTok{(}\AttributeTok{title =} \FunctionTok{paste}\NormalTok{(}\StringTok{"🌍 Réseau commercial mondial {-}"}\NormalTok{, annee),}
         \AttributeTok{subtitle =} \StringTok{"Graphique du réseau dirigé non pondéré (exportations)"}\NormalTok{) }\SpecialCharTok{+}
    \FunctionTok{theme\_void}\NormalTok{()}
\NormalTok{\}}

\CommentTok{\# 🔍 Visualisation des graphes pour 1900, 1955, 2009}
\FunctionTok{creer\_graphe\_reseau}\NormalTok{(trade\_clean, }\DecValTok{1900}\NormalTok{)}
\end{Highlighting}
\end{Shaded}

\pandocbounded{\includegraphics[keepaspectratio]{Examen_visualisation_trade-networks_files/figure-latex/unnamed-chunk-14-1.pdf}}

\begin{Shaded}
\begin{Highlighting}[]
\FunctionTok{creer\_graphe\_reseau}\NormalTok{(trade\_clean, }\DecValTok{1955}\NormalTok{)}
\end{Highlighting}
\end{Shaded}

\begin{verbatim}
## Warning: ggrepel: 21 unlabeled data points (too many overlaps). Consider
## increasing max.overlaps
\end{verbatim}

\pandocbounded{\includegraphics[keepaspectratio]{Examen_visualisation_trade-networks_files/figure-latex/unnamed-chunk-14-2.pdf}}

\begin{Shaded}
\begin{Highlighting}[]
\FunctionTok{creer\_graphe\_reseau}\NormalTok{(trade\_clean, }\DecValTok{2009}\NormalTok{)}
\end{Highlighting}
\end{Shaded}

\begin{verbatim}
## Warning: ggrepel: 166 unlabeled data points (too many overlaps). Consider
## increasing max.overlaps
\end{verbatim}

\pandocbounded{\includegraphics[keepaspectratio]{Examen_visualisation_trade-networks_files/figure-latex/unnamed-chunk-14-3.pdf}}

`

\end{document}
